\begin{abstract}
Understanding how neural networks implement learned capabilities requires moving beyond individual neurons to analyzing circuit-level structure. We present a systematic graph-theoretic analysis of feedforward loop (FFL) motifs in 200 neural network attribution graphs from the Neuronpedia platform, spanning 8 capability domains (antonym, arithmetic, code completion, country--capital, multi-hop reasoning, rhyme, sentiment, translation). FFL motifs are universally over-represented (median $Z=47.2$, mean $Z=46.2$) across all 8 domains, with per-domain Z-scores ranging from 26.3 to 59.4. Node ablation experiments on 11,185,566 FFL instances identify 20,056 hub nodes whose removal causes 1.92$\times$ greater downstream attribution loss than layer-matched controls, with a dose--response correlation of $r=0.88$ ($p<0.001$). Weighted motif features (Onnela intensity, sign coherence, path dominance) achieve spectral clustering $NMI=0.705$ against domain labels, vastly outperforming binary motif counts ($NMI=0.101$) and matching graph-level statistics ($NMI=0.677$); combining all features yields $NMI=0.844$ ($p=0.001$). Unique information decomposition confirms motifs carry genuinely orthogonal structural information (unique $R^2=0.0184$, CI excludes 0; residualized $NMI=0.2643$, $p=0.001$; MI $retained=23.6\%$). All 4 universal 4-node motif types (IDs 77, 80, 82, 83) show 100\% FFL containment, confirming higher-order universality is fully FFL-derivative. These results establish FFL motifs as fundamental organizational primitives of neural network attribution circuits (Table~\ref{table:corpus}; Table~\ref{table:zscores}) \cite{conmy2023automated, marks2024sparse}.
\end{abstract}